\SourceLeaf{242}Fläche; $V'-\varepsilon U$ in der Fläche, so weit sie bei der zweiten Vertheilung belegt ist, und daher in demselben Stücke der Fläche auch $v-U$, weil $v=\frac{V'-V^0}{\varepsilon}$.

Ob in der zweiten Vertheilung die ganze Fläche belegt ist, oder ob ein grösseres oder kleineres Stück unbelegt bleibt, wird von dem Coëfficienten $\varepsilon$ abhangen. Da die zweite Vertheilung in die erste übergeht, wenn $\varepsilon=0$ wird, so wird allgemein zu reden das für einen bestimmten Werth von $\varepsilon$ unbelegt gebliebene Stück der Fläche sich verengern, wenn $\varepsilon$ abnimmt, und sich schon ganz füllen, ehe $\varepsilon$ den Werth $0$ erreicht hat. In singulären Fällen aber kann es sich auch so verhalten, dass immer ein Stück unbelegt bleibt, so lange $\varepsilon$ von $0$ verschieden ist und nicht das entgegengesetzte Zeichen annimmt. Für unsern Zweck ist es zureichend, $\varepsilon$ unendlich klein anzunehmen, wo sich leicht nachweisen lässt, dass jedenfalls kein endliches Flächenstück unbelegt bleiben kann. Denn im entgegengesetzten Falle würde nach der Schlussbemerkung des Art.\ 32 das Integral $\int V'm'\mathrm ds$ um einen endlichen Unterschied grösser sein müssen als $\int V^0m^0\mathrm ds$: wird dieser Unterschied mit $e$ bezeichnet, so ist der Unterschied der beiden Integrale
\[
\int(V'-2\varepsilon U)m'\mathrm ds-\int(V^0-2\varepsilon U)m^0\mathrm ds=e-2\varepsilon\int U(m'-m^0)\mathrm ds
\]
welcher für ein unendlich kleines $\varepsilon$ einen positiven Werth behält, im Widerspruch mit der Voraussetzung, dass $\int(V-2\varepsilon U)m\mathrm ds$ in der zweiten Vertheilung seinen Minimumwerth hat.

Man schliesst hieraus, dass wenn man in der dritten Vertheilung für $\mu$ den Grenzwerth von $\frac{m'-m^0}{\varepsilon}$, bei unendlicher Abnahme von $\varepsilon$, annimmt, $v-U$ in der ganzen Fläche einen constanten Werth hat.

Bilden wir nun eine vierte Vertheilung, wobei $m=m^0+\mu$ gesetzt wird, die ganze Masse also $=M$ bleibt, so wird das daraus entspringende Potential $=V^0+v$ sein, mithin in der ganzen Fläche die Grösse $U$ um die constante Differenz $V^0+v-U$ übertreffen, wodurch also der oben ausgesprochene Lehrsatz erwiesen ist.

\Article{34}
Es bleibt noch übrig, zu beweisen, dass nur Eine Vertheilungsart einer gegebenen Masse $M$ möglich ist, bei welcher $V-U$ in der ganzen Fläche constant ist. In der That, gäbe es zwei verschiedene dies leistende Vertheilungsar\SourceLeaf{243}ten, so würde, wenn man $m$ und $V$ in der ersten mit $m'$, $V'$, in der zweiten mit $m''$, $V''$ bezeichnet, von einer dritten Massenvertheilung, in welcher $m=m'-m''$ angenommen wird, das Potential $=V'-V''$ und folglich constant sein, und die Gesammtmasse $=0$. Das constante Potential müsste daher nach Art.\ 27 nothwendig $=0$ sein, und folglich nach Art.\ 28 auch $m'-m''=0$, oder die beiden Vertheilungen identisch.

Endlich muss noch erwähnt werden, dass es immer eine Massenvertheilung gibt, wobei die Differenz $V-U$ einen \textit{gegebenen} constanten Werth erhält. Bedeutet nemlich $\alpha$ einen beliebigen constanten Coëfficienten, so wird, indem wir die Bezeichnungen für die erste und dritte Vertheilung im vorhergehenden Artikel beibehalten, das Potential derjenigen Vertheilung, wobei $m=\alpha m^0+\mu$ angenommen wird, $=\alpha V^0+v$ sein, und dem constanten Unterschiede $\alpha V^0+v-U$ durch gehörige Bestimmung des Coëfficienten $\alpha$ jeder beliebige Werth ertheilt werden können. Die Gesammtmasse dieser Vertheilung ist dann aber nicht mehr willkürlich, sondern $=\alpha M$. Übrigens erhellt auf dieselbe Art wie vorhin, dass auch diese Vertheilungsbedingung nur auf eine einzige Art erfüllt werden kann.

\Article{35}
Die wirkliche Bestimmung der Vertheilung der Masse auf einer gegebenen Fläche für jede vorgeschriebene Form von $U$ übersteigt in den meisten Fällen die Kräfte der Analyse in ihrem gegenwärtigen Zustande. Der einfachste Fall, wo sie in unsrer Gewalt ist, ist der einer ganzen Kugelfläche; wir wollen jedoch sofort den allgemeinern behandeln, wo die Fläche von der Kugelfläche sehr wenig abweicht, und Grössen von höherer Ordnung, als die Abweichung selbst, vernachlässigt werden dürfen.

Es sei $R$ der Halbmesser der Kugel, $r$ die Entfernung jedes Punktes im Raume von ihrem Mittelpunkte, $u$ der Winkel zwischen $r$ und einer festen geraden Linie, $\lambda$ der Winkel zwischen der durch diese gerade Linie und $r$ gelegten Ebene und einer festen Ebene. Der Abstand eines unbestimmten Punktes in der gegebenen geschlossenen Fläche vom Mittelpunkte der Kugel sei $=R(1+\gamma z)$, wo $\gamma$ ein constanter sehr kleiner Factor ist, dessen höhere Potenzen vernachlässigt werden, $z$ hingegen eben so wie $U$ Functionen von $u$ und $\lambda$.

\SourceLeaf{244}Das Potential $V$ der auf die Kugelfläche vertheilten Masse wird in jedem Punkte des äussern Raumes durch eine nach Potenzen von $r$ fallende Reihe ausgedrückt werden, welcher wir die Form geben
\[
A^0\frac Rr+A'\left(\frac Rr\right)^2+A''\left(\frac Rr\right)^3+\text{u.\,s.\,f.}
\]
in jedem Punkte des innern Raumes hingegen durch die steigende Reihe
\[
B^0+B'\frac rR+B''\left(\frac rR\right)^2+B'''\left(\frac rR\right)^3+\text{u.\,s.\,f.}
\]
Die Coëfficienten $A^0,A',A''$ u.\,s.\,f.\ sind Functionen von $u$ und $\lambda$, welche bekannten partiellen Differentialgleichungen Genüge leisten, \textit{Allg.\ Th.\ d.\ Erdm.}\ Art.\ 18, und eben so $B^0,B',B''$ u.\,s.\,f.\ Auf der vorgegebenen Fläche soll nun das Potential einer gegebenen Function von $u$ und $\lambda$ gleich werden, nemlich $V=U$, also
\[
\left(\frac rR\right)^{\frac12}V=(1+\gamma z)^{\frac12}U
\]
Nehmen wir also an, dass $(1+\gamma z)^{\frac12}U$ in eine Reihe
\[
P^0+P'+P''+P'''+\text{u.\,s.\,w.}
\]
entwickelt sei, dergestalt, dass die einzelnen Glieder $P^0$, $P'$, $P''$, $P'''$ u.\,s.\,f.\ gleichfalls den gedachten Differentialgleichungen Genüge leisten, und erwägen, dass die beiden obigen Reihen für das Potential bis zur Fläche selbst gültig bleiben müssen, so erhellt, dass
\[
\begin{aligned}
&P^0+P'+P''+P'''+\text{u.\,s.\,f.}\\
&\qquad=A^0(1+\gamma z)^{-\frac12}+A'(1+\gamma z)^{-\frac32}+A''(1+\gamma z)^{-\frac52}+\text{u.\,s.\,f.}\\
&\qquad=B^0(1+\gamma z)^{\frac12}+B'(1+\gamma z)^{\frac32}+B''(1+\gamma z)^{\frac52}+\text{u.\,s.\,f.}
\end{aligned}
\]
sein wird. Wir schliessen hieraus, dass, wenn man Grössen der Ordnung $\gamma$ vernachlässigt,
\[
P^0+P'+P''+\text{u.\,s.\,f.}=A^0+A'+A''+\text{u.\,s.\,f.}
\]
und also (da eine Function von $u$, $\lambda$ nur auf Eine Art in eine Reihe entwickelt werden kann, deren Glieder den erwähnten Differentialgleichungen Genüge leisten)
\[
P^0=A^0,\qquad P'=A',\qquad P''=A''\quad\text{u.\,s.\,f.}
\]
Eben so wird, Grössen der Ordnung $\gamma$ vernachlässigt,
\[
P^0=B^0,\qquad P'=B',\qquad P''=B''\quad\text{u.\,s.\,f.}
\]

\SourceLeaf{245}Setzt man also (I)
\[
\begin{aligned}
A^0&=P^0+\gamma a^0,& B^0&=P^0-\gamma b^0\\
A'&=P'+\gamma a',& B'&=P'-\gamma b'\\
A''&=P''+\gamma a'',& B''&=P''-\gamma b''\\
A'''&=P'''+\gamma a''',& B'''&=P'''-\gamma b'''\\
&\text{u.\,s.\,f.}
\end{aligned}
\]
wo offenbar auch $a^0,a',a'',a'''$ u.\,s.\,f., ingleichen $b^0,b',b'',b'''$ u.\,s.\,f.\ den erwähnten Differentialgleichungen Genüge leisten werden, und substituirt diese Werthe in den obigen Gleichungen, indem man dabei Grössen von der Ordnung $\gamma\gamma$ vernachlässigt, so wird, nachdem mit $\gamma$ dividirt ist, bis auf Fehler von der Ordnung $\gamma$ genau
\[
\begin{aligned}
a^0+a'+a''+a'''+\text{u.\,s.\,f.}&=\tfrac12 z(P^0+3P'+5P''+7P'''+\text{u.\,s.\,f.})\\
b^0+b'+b''+b'''+\text{u.\,s.\,f.}&=\tfrac12 z(P^0+3P'+5P''+7P'''+\text{u.\,s.\,f.})
\end{aligned}
\]
Es ist also, bis auf Fehler der Ordnung $\gamma$ genau,
\[
b^0=a^0,\qquad b'=a',\qquad b''=a''\quad\text{u.\,s.\,f.}
\]
und folglich, bis auf Fehler der Ordnung $\gamma\gamma$ genau, (II)
\[
B^0=P^0-\gamma a^0,\qquad B'=P'-\gamma a',\qquad B''=P''-\gamma a''\quad\text{u.\,s.\,f.}
\]

Der Differentialquotient $\frac{\mathrm dV}{\mathrm dr}$ hat in der Fläche selbst zwei verschiedene Werthe, und der auf ein negatives $\mathrm dr$ oder auf die innere Seite sich beziehende übertrifft den auf der äussern Seite geltenden um $4\pi m\cos\theta$, wenn $m$ die Dichtigkeit an der Durchschnittsstelle und $\theta$ den Winkel zwischen $r$ und der Normale bezeichnet (Art.\ 13, wo $t,A,k^0$ dasselbe bedeuten, was hier $r,\theta,m$ sind). Man findet diese beiden Werthe, wenn man die beiden im innern und äussern Raume geltenden Ausdrücke für $V$ nach $r$ differentiirt, und dann $r=R(1+\gamma z)$ setzt. Es ist also der erste
\[
=\frac1R\{B'+2B''(1+\gamma z)+3B'''(1+\gamma z)^2+\text{u.\,s.\,f.}\}
\]
und der zweite
\[
-\frac1R\{A^0(1+\gamma z)^{-2}+2A'(1+\gamma z)^{-3}+3A''(1+\gamma z)^{-4}+\text{u.\,s.\,f.}\}
\]

\SourceLeaf{246}Wir haben also, wenn wir die Differenz mit $R(1+\gamma z)^{\frac32}$ multipliciren,
\[
\begin{aligned}
&4\pi mR\cos\theta.(1+\gamma z)^{\frac32}\\
&\quad=\phantom{{}+{}}A^0(1+\gamma z)^{-\frac12}+2A'(1+\gamma z)^{-\frac32}+3A''(1+\gamma z)^{-\frac52}+\text{u.\,s.\,f.}\\
&\qquad+B'(1+\gamma z)^{\frac32}+2B''(1+\gamma z)^{\frac52}+3B'''(1+\gamma z)^{\frac72}+\text{u.\,s.\,f.}
\end{aligned}
\]
Substituiren wir hierin statt $A^0,A'$ u.\,s.\,f.\ die Werthe aus I, und statt $B^0,B'$ u.\,s.\,w.\ die Werthe aus II, und lassen weg, was von der Ordnung $\gamma\gamma$ ist, so erhalten wir
\[
\begin{aligned}
4\pi mR\cos\theta.(1+\gamma z)^{\frac32}&=P^0+3P'+5P''+7P'''+\text{u.\,s.\,f.}\\
&\quad+\gamma(a^0+a'+a''+a'''+\text{u.\,s.\,f.})\\
&\quad-\tfrac12\gamma z(P^0+3P'+5P''+\text{u.\,s.\,f.})
\end{aligned}
\]
folglich, da die beiden letzten Reihen bis auf Grössen der Ordnung $\gamma\gamma$ einander destruiren,
\[
m=\frac{(1+\gamma z)^{-\frac32}}{4\pi R\cos\theta}.(P^0+3P'+5P''+7P'''+\text{u.\,s.\,f.})
\]
womit die Aufgabe gelöst ist. Anstatt $(1+\gamma z)^{-\frac32}$ kann man auch schreiben $1-\frac32\gamma z$, und den Divisor $\cos\theta$ weglassen, insofern, wenigstens allgemein zu reden, $\theta$ von der Ordnung $\gamma$, und also $\cos\theta$ von $1$ nur um eine Grösse der Ordnung $\gamma\gamma$ verschieden ist.

Für den Fall einer Kugel, wo $\gamma=0$, hat man in aller Schärfe
\[
m=\tfrac1{4\pi R}(P^0+3P'+5P''+7P'''+\text{u.\,s.\,f.})
\]
indem $P^0+P'+P''+P'''+\text{u.\,s.\,f.}$ die Entwicklung von $U$ selbst vorstellt.

\Article{36}
Die Grösse $U$ ist in den bisherigen Untersuchungen unbestimmt gelassen: die Anwendung derselben auf den Fall, wo für $U$ das Potential eines gegebenen Massensystems angenommen wird, bahnt uns nun den Weg zu folgendem wichtigen

\textsc{Lehrsatz.} Anstatt einer beliebigen gegebenen Massenvertheilung $D$, welche entweder bloss auf den innern von einer geschlossenen Fläche $S$ begrenzten Raum beschränkt ist, oder bloss auf den äussern Raum, lässt sich eine Massenvertheilung $E$ bloss auf der Fläche selbst substituiren, mit dem Erfolge, dass die
\SourceLeaf{247}Wirkung von $E$ der Wirkung von $D$ gleich wird, in allen Punkten des äussern Raumes für den ersten Fall, oder in allen Punkten des innern Raumes für den zweiten.

Es wird dazu nur erfordert, dass, indem das Potential von $D$ in jedem Punkte von $S$ mit $U$, das Potential von $E$ hingegen mit $V$ bezeichnet wird, in der ganzen Fläche für den ersten Fall $V-U=0$, für den zweiten aber nur constant werde. Es wird nemlich $-U$ das Potential einer Vertheilung $D'$ sein, die der $D$ entgegengesetzt ist (so dass an die Stelle jedes Massentheils ein entgegengesetztes tritt), also $V-U$ das Potential der zugleich bestehenden Vertheilungen $D'$ und $E$; die Wirkungen daraus werden sich folglich im ersten Fall im ganzen äussern Raume, im zweiten im ganzen innern destruiren (Artt.\ 27 und 25), oder die Wirkungen von $D$ und $E$ werden in den betreffenden Räumen gleich sein. Übrigens wird die ganze Masse in $E$ für den ersten Fall der Masse in $D$ gleich sein, im zweiten aber willkürlich bleiben.

Der Lehrsatz, welcher in der \textit{Intensitas vis magneticae} Art.\ 2 angekündigt, und auch in der \textit{Allgemeinen Theorie des Erdmagnetismus} an verschiedenen Stellen angeführt ist, erscheint jetzt als ein specieller Fall des hier bewiesenen.

\Article{37}
Obgleich, wie schon im 35.\ Artikel bemerkt ist, die wirkliche vollständige Ausmittelung der Vertheilung $E$ in den meisten Fällen unüberwindliche Schwierigkeiten darbietet, so gibt es doch einen, wo sie mit grosser Leichtigkeit geschehen kann, und der hier noch besonders angeführt zu werden verdient. Dies ist nemlich der, wo $U$ constant, also $S$ eine Gleichgewichtsfläche für das gegebene Massensystem $D$ ist. Man sieht leicht, dass hier nur von dem Falle die Rede zu sein braucht, wo $D$ im innern Raume angenommen wird, und nicht die Gesammtmasse $=0$ ist, da sonst gar keine Wirkung da sein würde, die durch eine Massenvertheilung auf $S$ ersetzt zu werden brauchte.

Es sei $O$ ein Punkt der Fläche $S$, und $r$ eine gerade Linie, welche die Fläche daselbst unter rechten Winkeln schneidet, und in der Richtung von Innen nach Aussen als wachsend betrachtet wird; es sei ferner $-C$ der Werth des Differentialquotienten $\frac{\mathrm dU}{\mathrm dr}$ in $O$, und $m$ die Dichtigkeit, welche bei der Massenvertheilung $E$ in $O$ Statt hat. Der Differentialquotient $\frac{\mathrm dV}{\mathrm dr}$ wird in $O$ zwei verschiedene Werthe haben; der auf die äussere Seite sich beziehende wird,
\SourceLeaf{248}weil in der Fläche und im ganzen äussern Raume $V=U$ ist, dem Differentialquotienten $\frac{\mathrm dU}{\mathrm dr}$ gleich, also $=-C$ sein; der auf die innere Seite sich beziehende hingegen $=0$, weil $V$ in der Fläche und im ganzen innern Raume constant ist. Da nun aber der zweite Werth um $4\pi m$ grösser ist als der erste, so haben wir $4\pi m=C$ oder $m=\frac C{4\pi}$. Offenbar ist $C$ nichts anderes, als die aus der Massenvertheilung $D$ entspringende Kraft, und hat mit der Gesammtmasse einerlei Zeichen.

\clearpage
\thispagestyle{empty}\SourceLeaf{249}
\pdfbookmark[0]{Dioptrische Untersuchungen}{dioptrische-title}
\begin{center}
\vspace*{0.6in}
{\Large DIOPTRISCHE UNTERSUCHUNGEN\par}
\vspace{0.6in}
VON\par
\vspace{0.5in}
{\large CARL FRIEDRICH GAUSS\par}
\vspace{0.4in}
DER KÖNIGL. SOCIETÄT ÜBERGEBEN MDCCCXL DECEMBER X.\par
\vfill
\rule{0.85\textwidth}{0.3pt}\par
Abhandlungen der königlichen Gesellschaft der Wissenschaften zu Göttingen Band I.\par
Göttingen \textsc{mdcccxliii}.\par
\rule{0.85\textwidth}{0.3pt}
\end{center}
\clearpage
\SourceLeaf{251}
\fancyhead[L]{\small Dioptrische Untersuchungen}
\renewcommand{\Article}[1]{\par\medskip\Needspace{4\baselineskip}\pdfbookmark[1]{Art. #1}{dioptrische-art-#1}\centerline{\textbf{#1.}}\nobreak\smallskip}
\pdfbookmark[0]{Dioptrische Untersuchungen}{dioptrische-body}
\begin{center}
DIOPTRISCHE UNTERSUCHUNGEN.
\end{center}
\medskip
Die Betrachtung des Weges, welchen durch Linsengläser solche Lichtstrahlen nehmen, die gegen die gemeinschaftliche Axe derselben sehr wenig geneigt sind, und der davon abhangenden Erscheinungen, bietet sehr elegante Resultate dar, welche durch die Arbeiten von \textsc{Cotes}, \textsc{Euler}, \textsc{Lagrange} und \textsc{Möbius} erschöpft scheinen könnten, aber doch noch mehreres zu wünschen übrig lassen. Ein wesentlicher Mangel der von jenen Mathematikern aufgestellten Sätze ist, dass dabei die Dicke der Linsen vernachlässigt wird, wodurch ihnen ein ihren Werth sehr verringernder Charakter von Ungenauigkeit und Naturwidrigkeit aufgeprägt wird. Ohne in Abrede zu stellen, dass für manche andere dioptrische Untersuchungen, namentlich für diejenigen, wobei die sogenannte Abweichung wegen der Kugelgestalt der Linsenflächen in Betracht gezogen wird, die anfängliche Vernachlässigung der Dicke der Linsen sehr nützlich, ja nothwendig wird, um einfachere und geschmeidigere Vorschriften für Überschläge und erste Annäherungen zu gewinnen, wird man sich doch gern einer solchen Aufopferung aller Schärfe da enthoben sehen, wo es ohne allen oder ohne erheblichen Verlust für die Einfachheit der Resultate geschehen kann. Auf einen den mathematischen Sinn unangenehm berührenden Mangel an Präcision stossen wir zum Theil schon bei den ersten Begriffsbestimmungen der Dioptrik. Die Begriffe von Axe und Brennpunkt einer Linse stehen zwar mit Schärfe fest; allein nicht so ist es mit
